\chapter{Der Offline-Test}

Erst dieser Test zeigt, ob Ihr System im Notfall wirklich ohne Internet funktioniert. Machen Sie ihn \textbf{einmal gründlich, solange das Internet noch erreichbar ist}, damit Sie fehlende Inhalte nachladen können. Wiederholen Sie ihn nach jeder größeren Änderung (Kapitel~9: Tabelle „Offline-Test“).

\section{Vorbereitung}
\begin{itemize}[leftmargin=1.5em,itemsep=3pt,label={\abhaken}]
\item Alle Downloads sind fertig (Kapitel~5, Seite „Downloads“).
\item Sie haben die Notfall-Karte (Kapitel~7) zur Hand.
\end{itemize}

\section{Internet trennen}
\begin{enumerate}[leftmargin=1.5em,itemsep=3pt]
\item \textbf{Netzwerkkabel ziehen.} Bei WLAN: oben rechts im Desktop auf das Netzwerksymbol klicken und WLAN ausschalten.
\item \textbf{Beweis:} Öffnen Sie das Terminal (\taste{Strg}+\taste{Alt}+\taste{T}) und tippen Sie:
\befehl{ping -c 2 8.8.8.8}
Es muss „100\,\% packet loss“ erscheinen (Abbildung~\ref{fig:offline}). Sonst sind Sie noch im Internet.
\end{enumerate}
\bild[0.8]{44-offline-test-terminal}{Offline-Beweis im Terminal (hier zusätzlich mit der Liste der geladenen Inhaltsdateien).}\label{fig:offline}

\section{Neu starten ohne Internet}
Starten Sie den Rechner \textbf{neu, ohne Netzwerk}. Das ist der eigentliche Test: Im Notfall startet er genau so. Warten Sie etwa zwei Minuten nach dem Desktop.

\section{Die Checkliste}
\renewcommand{\arraystretch}{1.5}
\begin{center}
\begin{tabularx}{\linewidth}{@{}c>{\RaggedRight\arraybackslash}X>{\RaggedRight\arraybackslash}p{3.2cm}@{}}
\toprule
 & \textbf{Prüfung} & \textbf{Ergebnis} \\
\midrule
\abhaken & Der Rechner startet ohne Netzwerkkabel bis zum Desktop. & \\
\abhaken & Im Browser \texttt{http://localhost:8080} zeigt die Kommandozentrale. & \\
\abhaken & Unter \menue{Docs} lassen sich die deutschen Anleitungen lesen. & \\
\abhaken & Die \menue{Wissensbibliothek} (Kiwix, \texttt{http://localhost:8090}) zeigt alle geladenen Bücher. & \\
\abhaken & Eine Suche nach einem Begriff liefert einen Artikel (zum Beispiel „Wasser“ oder „water“). & \\
\abhaken & Wikipedia öffnet einen Artikel samt Seitenlinks. & \\
\abhaken & (falls geladen) Karten zeigen den gewählten Ausschnitt. & \\
\abhaken & (falls eingerichtet) Die KI antwortet auf eine Frage. & \\
\bottomrule
\end{tabularx}
\end{center}

\bild[0.85]{45-kiwix-offline}{Die Wissensbibliothek (Kiwix) ohne Internet: Alle geladenen Bücher sind da.}
\bild[0.85]{47-wikipedia-offline}{Ein Wikipedia-Artikel ohne Internetverbindung.}

\hinweis{Zeigt NOMAD Bilder oder Texte in \emph{anderer} Sprache als Deutsch, ist das richtig: Die geladenen Bücher sind überwiegend englisch (Anhang).}

\section{Wenn etwas nicht geht}
\begin{itemize}[leftmargin=1.5em,itemsep=3pt]
\item \textbf{Seite nicht erreichbar:} noch eine Minute warten, dann Notfall-Karte (Kapitel~7), Punkt „NOMAD starten“.
\item \textbf{Ein Buch fehlt:} Internet wieder anschließen, Download im Content-Manager wiederholen, Test erneut machen.
\item \textbf{Zeigt nur eine weiße Seite:} Browser mit \taste{Strg}+\taste{F5} neu laden.
\end{itemize}
